\section{Experimental Methodology}
\label{sec:method}

\setlength\heavyrulewidth{1pt}
\begin{table}[t]
  \centering
  \caption{NVMe and NVM server properties.}
  \resizebox{0.7\columnwidth}{!}{%
    \begin{tabular}{llclc}
      \multicolumn{1}{c}{\textbf{ID}} & \multicolumn{1}{c}{\textbf{CPU}}  & \multicolumn{1}{c}{\textbf{\begin{tabular}[c]{@{}c@{}}DRAM\end{tabular}}} & \multicolumn{1}{c}{\textbf{Storage}}                                                & \multicolumn{1}{c}{\textbf{Kernel}} \\
        \toprule
        \begin{tabular}[c]{@{}l@{}}NVMe \\ Server\end{tabular}  & \begin{tabular}[c]{@{}l@{}}Intel Xeon E5-2630 \\ 32 Cores @ 2.4 GHz\end{tabular} & \SI{256}{GB} & \begin{tabular}[c]{@{}l@{}}\SI{2}{TB} Samsung PM983 \\ PCI Express NVMe SSD\end{tabular} & 4.14 \\ \hline
        \begin{tabular}[c]{@{}l@{}}SATA \\ Server\end{tabular}  & \begin{tabular}[c]{@{}l@{}}Intel Xeon E5-2630 \\ 32 Cores @ 2.4 GHz\end{tabular} & \SI{256}{GB} & \begin{tabular}[c]{@{}l@{}}$2 \times$ \SI{256}{GB} Samsung \\ SATA SSD 850\end{tabular} & 4.14 \\ \hline
        \begin{tabular}[c]{@{}l@{}}NVM \\ Server\end{tabular} & \begin{tabular}[c]{@{}l@{}}Intel Xeon Platinum \\ 24 cores @ 2.4 GHz\end{tabular} & \SI{192}{GB} & \begin{tabular}[c]{@{}l@{}}\SI{3}{TB} Intel Optane DC \\ Persistent Memory\end{tabular}  & 3.10 \\ \bottomrule

    \end{tabular}%
    }
    \label{tab:server}
\end{table}

We answer the following questions in our evaluation:
\begin{enumerate}
  \item How does \teraheapName{} perform compared to native JVM and
    state-of-the-art Panthera with NVMe SSDs and NVM?
  \item What are the space requirements of H2 in the storage device
    and for its metadata in DRAM?
  \item What is the overhead of tracking references and moving
    objects to H2 during \teraheapGC?
  \item How does \teraheapName{} scale with an increasing number of mutator
    threads and dataset size?
  \item How does the performance of \teraheapName{} vary based on (1) the
    distribution of DRAM between H1 and the page cache for H2, and (2) the
    type of storage device used?
  \item How does \teraheapName{} perform compared to newer garbage
    collectors?
\end{enumerate}

\paragraph{\textbf{Server infrastructure}}
We evaluate \teraheapName{} both with block-addressable NVMe SSDs,
SATA SSDs, and byte-addressable NVM as the backing device
for H2. Table~\ref{tab:server} shows the properties of each server.
Our NVM server operates in two modes:
%
(1) \emph{App Direct mode} uses \SI{192}{GB} DRAM as main memory and
\SI{2}{TB} NVM as persistent storage device.
%
(2) \emph{Mixed mode} partitions NVM to use \SI{1}{TB} in memory mode
and \SI{2}{TB} in App Direct mode. DRAM (\SI{192}{GB}) acts as a cache
for \SI{1}{TB} NVM (memory mode) controlled by the CPU's memory
controller. In App Direct mode, the system mounts NVM on an ext4-DAX
file system to establish direct mappings to the device.

\paragraph{\textbf{Baseline and \teraheapName{} configurations}}
We use Spark v3.3.0 with the Kryo serializer~\cite{kryo}, a
state-of-the-art highly optimized \teraheapSD library for Java that Spark
recommends.
%
We run Spark with OpenJDK17. We run Giraph v1.2 with Hadoop v2.4
and OpenJDK8, as it does not support more recent versions of OpenJDK.
%
We use an executor with eight mutator threads for
both Spark and Giraph.
%
Also, we run GDS using Neo4j v.5.15 and GDS v.2.7 with OpenJDK17. For GDS, we
use four mutator threads, the maximum number of threads that the GDS community
edition supports~\cite{gds}.
%
Our experiments use two garbage collectors in different configurations: PS in
OpenJDK8 and OpenJDK17 and Garbage First (G1)~\cite{zhao:g1} in OpenJDK17. For
PS, we use 16 \teraheapGC threads.
G1 uses two parameters:
%
(1) the number of parallel \teraheapGC threads, which we set to eight (max
value) and
%
(2) the ratio of concurrent to mutator threads, which we set to two as
the recommended configuration, is one-fourth of the parallel \teraheapGC
threads~\cite{g1:manual}.
%
In our experiments, we utilize eight cores while disabling the remaining cores
of the servers. For the performance scaling experiments
(Section~\S\ref{sec:scaling}), we use 4, 8, and 16 cores to match the number
of mutator threads.

\begin{table}[t]
\centering
\caption{Summary of baselines.}
\resizebox{0.7\columnwidth}{!}{%
\begin{tabular}{l|cccc}
\multicolumn{1}{c|}{\textbf{Baseline}} &
  \textbf{DRAM} &
  \textbf{\begin{tabular}[c]{@{}c@{}}NVMe SSD\end{tabular}} &
  \textbf{\begin{tabular}[c]{@{}c@{}}NVM\\ (App Direct mode)\end{tabular}} &
  \textbf{\begin{tabular}[c]{@{}c@{}}NVM \\ (Memory mode)\end{tabular}} \\
    \toprule
    Spark-SD   & Heap       & Off-heap & -        & -    \\
    Spark-SD   & Heap       & -        & Off-heap & -    \\
    Spark-MO   & -          & -        & -        & Heap \\
    Giraph-OOC & Heap       & Off-Heap & -        & -    \\
    GDS-Heap   & -          & Heap     & -        & -    \\
    GDS-OldGen & Young gen. & Old gen. & -        & -    \\
\bottomrule
\end{tabular}%
}
\label{tab:baselines}
\end{table}

Table~\ref{tab:baselines} summarizes the Spark, Giraph, and GDS configurations
we use as baselines.
%
The two Spark-SD configurations place executor memory (heap) in DRAM
and cache RDDs in the on-heap cache, up to 50\% of the total heap
size. Any remaining RDDs are serialized in the off-heap cache, over
either NVMe SSD (first line of Table~\ref{tab:baselines}) or NVM in
App Direct mode (second line of Table~\ref{tab:baselines}).
%
Spark-MO places executor memory (heap) over NVM in memory mode,
caching all RDDs on-heap.
%
Giraph-OOC places the heap in DRAM and offloads vertices, edges, and
messages off-heap to the NVMe SSD.
%
GDS-Heap allocates the managed heap on an NVMe SSD using
memory mapped I/O by setting the \emph{-XX:AllocateHeapAt} runtime
flag.
%
GDS-OldGen allocates only the old generation of the managed heap on an
NVMe SSD using memory mapping (setting the \emph{-XX:AllocateOldGenAt}
runtime flag) and the young generation in DRAM.

We configure \teraheapName{} to allocate H1 on DRAM and H2 over a file in NVMe
SSD or NVM via memory-mapped I/O (\teraheapMMIO). The file in both NVMe and
NVM servers is mapped to the JVM virtual address space where the
application can access the data with regular load/store
instructions~\cite{apapag:fmap}. Our experiments show that machine
learning (ML) workloads in Spark access the individual elements of
cached RDD partitions sequentially. For this reason, we configure
\teraheapName{} for Spark ML workloads to use huge pages (\SI{2}{MB}) in H2 to
reduce the frequency of page faults. Instead of the native
\emph{mmap}, we use HugeMap~\cite{jmal:hugemap} a custom, open source,
\teraheapMMIO path that enables huge pages for file-backed mappings.

\renewcommand{\arraystretch}{1.1}
\begin{table}[t]
	\centering
  \caption{Configuration for each workload on NVMe and NVM servers for
  Spark-SD, Spark-MO, and \teraheapName{}.}
	\resizebox{0.8\columnwidth}{!}{%
		\begin{tabular}{p{0.01cm}l|ccccc}
			                                                                 &
			\multicolumn{1}{c|}{}
      & \textbf{\begin{tabular}[c]{@{}c@{}}NVMe Server\end{tabular}}        &
        \textbf{\begin{tabular}[c]{@{}c@{}}Dataset\end{tabular}} &
          \textbf{\begin{tabular}[c]{@{}c@{}}Spark-SD\end{tabular}}
            &
            \textbf{\begin{tabular}[c]{@{}c@{}}Spark-MO\end{tabular}} & \textbf{\begin{tabular}[c]{@{}c@{}}TeraHeap\end{tabular}}      \\
			                                                                 & \multicolumn{1}{r|}{\textbf{GB}}                                       &
			\multicolumn{1}{l}{\textbf{DRAM}}                                & \multicolumn{1}{c}{\textbf{Size}}                                      & \multicolumn{1}{c}{\textbf{Heap}}                             & \multicolumn{1}{c}{\textbf{Heap}}                          & \multicolumn{1}{c}{\textbf{H1}}                                                  \\
			\toprule
      \multirow{5}{*}{\rotatebox[origin=c]{90}{\textbf{GraphX}}}
      & PageRank (PR)
      & 80
      & 32
      & 64
      & 1024
      & 64
      \\

      &\begin{tabular}[c]{@{}l@{}} Connected Components (CC)\end{tabular}
      & 84
      & 32
      & 68
      & 1024
      & 68
      \\

      & Shortest Path (SSSP)
      & 58
      & 32
      & 42
      & 650
      & 42 \\

      &
                                                                       SVDPlusPlus (SVD)
                                                                       &
                                                                       40
                                                                       &
                                                                       2
                                                                       &
                                                                       24
                                                                       &
                                                                       500
                                                                       &
                                                                       24
                                                                       \\

                                                                       &
                                                                       Triangle
                                                                       Counts
                                                                       (TR)
                                                                       &
                                                                       80
                                                                       &
                                                                       2
                                                                       &
                                                                       64
                                                                       &
                                                                       64
                                                                       &
                                                                       64
                                                                       \\
                                                                       \bottomrule
                                                                       \multirow{4}{*}[-2.5pt]{\rotatebox[origin=c]{90}{\textbf{MLlib}}}
                                                                       &
                                                                       Linear
                                                                       Regression
                                                                       (LR)
                                                                       &
                                                                       70
                                                                       &
                                                                       256
                                                                       &
                                                                       54
                                                                       &
                                                                       1084
                                                                       &
                                                                       54
                                                                       \\

			                                                                 & Logistic Regression (LgR)                                                    & 70                                                          & 256                                                            & 54                                                          & 1084              & 54 \\

                                                                       & \begin{tabular}[c]{@{}l@{}}Support Vector Machine (SVM)\end{tabular}
                                                                       & 48
                                                                       & 256
                                                                       & 32
                                                                       & 620
                                                                       & 32 \\

			                                                                 & \begin{tabular}[c]{@{}l@{}}Naive Bayes Classifier (BC)\end{tabular}
                                                                         & 98
                                                                         & 21
                                                                         & 82
                                                                         & 82
                                                                         & 82 \\
			\bottomrule
			\rotatebox[origin=c]{90}{\textbf{SQL}}
      & RDD-RL
      & 63
      & 16
      & 47
      & 96
      & 47 \\
			\bottomrule
		\end{tabular}%
	}
  \label{tab:spark_conf}
\end{table}

\begin{table}[t]
	\centering
  \caption{Giraph-OOC and \teraheapName{} configurations for each workload on
  our NVMe server.}
	\resizebox{0.7\columnwidth}{!}{%
		\begin{tabular}{l|cccccc}
			\multicolumn{1}{c|}{}            &
      \textbf{\begin{tabular}[c]{@{}c@{}}NVMe Server\end{tabular}} &
        \textbf{\begin{tabular}[c]{@{}c@{}}Dataset\end{tabular}} &
          \multicolumn{2}{c}{\textbf{Giraph-OOC}} & \multicolumn{2}{c}{\textbf{TeraHeap}}                                  \\
			\multicolumn{1}{r|}{\textbf{GB}} & \textbf{DRAM} & \textbf{Size}                                               & \textbf{Heap}                          & \textbf{DR2}                          & \textbf{H1} & \textbf{DR2} \\
			\toprule
			PageRank (PR)                    & 85                                                             & 31                                                          & 70                                      & 15                                  & 50              & 35        \\
      \begin{tabular}[c]{@{}l@{}}Community Detection \\ Label Propagation (CDLP) \end{tabular}          & 85                                                             & 31                                                          & 70                                      & 15                                  & 60              & 25        \\
        \begin{tabular}[c]{@{}l@{}}Weakly Connected \\ Components (WCC) \end{tabular}         & 85                                                             & 31                                                          & 70                                      & 15                                  & 60              & 25        \\
			Breadth-first Search (BFS)       & 65                                                             & 31                                                          & 48                                      & 17                                  & 35              & 30        \\
			Shortest Path (SSSP)             & 90                                                             & 31                                                          & 75                                      & 15                                  & 50              & 40        \\
			\bottomrule
		\end{tabular}%
	}
  \label{tab:giraph_conf}
\end{table}

\begin{table}[t]
  \centering
\caption{Parameters we use with each GDS-Neo4j configuration in our experiments.}
  \resizebox{0.8\columnwidth}{!}{%
    \begin{tabular}{l|ccccccc}
      & \textbf{NVMe Server} & \textbf{Dataset} & \textbf{GDS-Heap} &
      \multicolumn{2}{c}{\textbf{GDS-OldGen}} & \multicolumn{2}{c}{\textbf{TeraHeap}} \\
      \multicolumn{1}{r|}{\textbf{GB}} & \textbf{DRAM}        & \textbf{Size}    & \textbf{DR2}  & \textbf{YoungGen}   & \textbf{DR2}  & \textbf{H1}       & \textbf{DR2}      \\
      \toprule
      PageRank (PR)                                                                          & 14 & 70 & 14 & 4 & 10 & 11 & 3 \\
      \begin{tabular}[c]{@{}l@{}}Community Detection\\ Label Propagation (CDLP)\end{tabular} & 14 & 70 & 14 & 4 & 10 & 11 & 3 \\
        \begin{tabular}[c]{@{}l@{}}Weakly Connected\\ Components (WCC)\end{tabular}            & 30 & 70 & 30 & 6 & 24 & 22 & 8 \\
			\bottomrule

    \end{tabular}%
    }
    \label{tab:neo4j}
\end{table}


In Spark-SD, to capture the effect of large datasets and limited DRAM
capacity~\cite{chen:hybrid}, we use a small heap size that caches a
limited number of RDDs on-heap and the rest off-heap (Spark-SD column
in Table~\ref{tab:spark_conf}). In Spark-MO we find and use the
minimum heap size that fits all the cached data on-heap (Spark-MO
column in Table~\ref{tab:spark_conf}). In Giraph-OOC, we
experimentally find the minimum heap size for each workload
(Giraph-OOC Heap column in Table~\ref{tab:giraph_conf}). \teraheapName{} uses
the same amount of DRAM as Spark-SD and Giraph-OOC but divides it
between H1 (DR1) and system (DR2). The DRAM devoted to system use
(DR2) includes the Spark and Giraph drivers and the kernel page cache
for I/O.

In GDS-OldGen, by default, the JVM sets the maximum value of the young
generation to be 80\% of the available DRAM to balance the tradeoff
between \teraheapGC efficiency and heap size. However, setting the maximum
size of the young generation to 80\% can lead to longer pause times
during garbage collection. Thus, we investigate different values for
the maximum size of the young generation and provide the best
configuration for our experiments. Specifically, for PR and WCC, the
young generation is 30\% of DRAM, while for CDLP is 20\% of the DRAM.

For the division of DRAM, we explore H1 sizes between 50\% and 90\% of
DRAM capacity, and we report results with a configuration hand-tuned
for each workload. Table~\ref{tab:spark_conf} and
Table~\ref{tab:giraph_conf} show the H1 size of \teraheapName{} in each
workload, in Spark and Giraph, respectively. DR2 is always \SI{16}{GB}
for Spark, whereas Table~\ref{tab:giraph_conf} shows the DR2 size for
Giraph-OOC and \teraheapName{} in each Giraph workload. We limit the available
DRAM capacity in our experiments with the NVMe server using
\emph{cgroups}. Table~\ref{tab:spark_conf} and
Table~\ref{tab:giraph_conf} show the total DRAM capacity in the NVMe
server for each workload.

\paragraph{\textbf{Workloads and datasets}}
We use ten memory-intensive workloads from the Spark-Bench
suite~\cite{Li:SparkBench}. For Giraph and GDS, we use workloads from the LDBC
Graphalytics suite~\cite{iosup:ldbc}. SparkBench provides a data generator which
we use to generate input datasets for our Spark workloads. For Giraph and GDS
workloads, we use the datagen-9\_0-fb dataset~\cite{iosup:ldbc} and
datagen-9\_4-fb dataset~\cite{iosup:ldbc}, respectively.
Table~\ref{tab:spark_conf}, Table~\ref{tab:giraph_conf}, and
Table~\ref{tab:neo4j} show the dataset size for each workload.

\paragraph{\textbf{Execution time breakdowns and \teraheapSD overhead}}
Our experiments are long-running, and to navigate the design space in a
practical timeframe, we repeat each experiment five times, reporting the average
end-to-end execution time. We observe low run-to-run variation (up to 2\%), and
five iterations suffice. We break execution time into four components:
\emph{other} time, \teraheapSD + I/O time, minor \teraheapGC time, and major \teraheapGC time. Other
time includes mutator threads time. In \teraheapName{}, the other time potentially
includes I/O wait due to page faults when accessing the H2 backing device. In
Spark-SD (see Table~\ref{tab:baselines}), \teraheapSD time includes \teraheapSD time both for
shuffle and caching. In \teraheapName{} and Spark-MO (see Table~\ref{tab:baselines}),
all \teraheapSD time is due to shuffling. The JVM reports the time spent in minor and
major \teraheapGC.

To estimate \teraheapSD overhead, which occurs in mutator threads, we use a
sampling profiler~\cite{async-profiler} to collect execution samples
from the mutator threads. The samples include the stack trace, similar
to the flame graph~\cite{flamegraphs} approach. Then we group the
samples for all the paths that originate from the top-level
\texttt{writeObject()} and \texttt{readObject()} methods of the
KryoSerializationStream and KryoDeserializationStream classes. These
samples include both \teraheapSD for the compute cache and the shuffle
network path of Spark. We then use the ratio of \teraheapSD samples to the
total application thread samples as an estimate of the time spent in
\teraheapSD, and we plot this separately in our execution time breakdowns. We
run the profiler with a \SI{10}{ms} sampling interval, verifying that
this does not create significant overhead (less than 2\% of total
execution time).
